% Part: first-order-logic
% Chapter: axiomatic-deduction
% Section: provability-quantifiers

% verification of properties of provability needed for maximally
% consistent sets in the completeness chapter.

\documentclass[../../../include/open-logic-section]{subfiles}

\begin{document}

\olfileid{fol}{axd}{qpr}

\olsection{\usetoken{S}{derivability} والمكمّمات}

\begin{explain}
  ومن مقدمات برهان الاكتمال أن تثبت بالاشتقاقات البديهية هذه
  النتائج في علاقة~$\Proves$؛ فنبيّنها الآن.
\end{explain}

\begin{thm}
\ollabel{thm:strong-generalization} إذا كان ${\OLId{latin-ordinary}{c}}$ !!a{constant} لا يرد
في ${\OLId{greek-variable}{Gamma}}$ ولا في $!A({\OLId{latin-ordinary}{x}})$، وكان ${\OLId{greek-variable}{Gamma}} \Proves !A({\OLId{latin-ordinary}{c}})$، فإن ${\OLId{greek-variable}{Gamma}}
\Proves \lforall[{\OLId{latin-ordinary}{x}}][!A({\OLId{latin-ordinary}{x}})]$.
\end{thm}

\begin{proof}
% تصحيح برهاني (0123-I1): لا يلزم أن يكون c من مخزون الثوابت المميزة؛ نجدد ثوابت البرهان ثم نبدله بثابت من E.
خذ برهانًا منتهيًا~$\Pi$ لـ~$!A({\OLId{latin-ordinary}{c}})$ من~${\OLId{greek-variable}{Gamma}}$، وسمّ فروضه المستعملة
${\OLId{greek-variable}{Gamma}}_{\OLMathNumeral{0}}$. جدّد ثوابته المميَّزة بـ
  \olref[qua]{lem:qr-freshening} مع مجموعة المنع~$\{{\OLId{latin-ordinary}{c}}\}$، ثم اختر
${\OLId{latin-ordinary}{e}}\in\mathcal E\setminus\{{\OLId{latin-ordinary}{c}}\}$ لا يقع في شيء من صيغ البرهان المجدَّد ولا يساوي ثابتًا
مميَّزًا مكتوبًا على عقدة~\QR{} فيه. بدّل~${\OLId{latin-ordinary}{c}}$ و~${\OLId{latin-ordinary}{e}}$ تقابليًا في البرهان
كله. ولا يتغير~${\OLId{greek-variable}{Gamma}}_{\OLMathNumeral{0}}$، إذ كلا الثابتين غير وارد فيه؛ كما تحفظ إعادة
التسمية البديهيات وإثبات المقدَّم وصورتي~\QR{} وشروطهما. فينتج
${\OLId{greek-variable}{Gamma}}_{\OLMathNumeral{0}}\Proves !A({\OLId{latin-ordinary}{e}})$.

  ونضم إلى ذلك الصيغة~$!A({\OLId{latin-ordinary}{e}})\lif(\ltrue\lif !A({\OLId{latin-ordinary}{e}}))$ من بديهية اللزوم
  الأولى، فنستخرج بإثبات المقدَّم
  ${\OLId{greek-variable}{Gamma}}_{\OLMathNumeral{0}}\Proves\ltrue\lif !A({\OLId{latin-ordinary}{e}})$. ولما كان~${\OLId{latin-ordinary}{e}}\in\mathcal E$ غير وارد
في~${\OLId{greek-variable}{Gamma}}_{\OLMathNumeral{0}}$ ولا~$\ltrue$ ولا~$!A({\OLId{latin-ordinary}{x}})$، جاز إجراء~\QR{}، فحصلنا على
${\OLId{greek-variable}{Gamma}}_{\OLMathNumeral{0}}\Proves\ltrue\lif\lforall[{\OLId{latin-ordinary}{x}}][!A({\OLId{latin-ordinary}{x}})]$. ونضم إليه
${\OLId{greek-variable}{Gamma}}_{\OLMathNumeral{0}}\Proves\ltrue$، فنستنتج بإثبات المقدَّم
${\OLId{greek-variable}{Gamma}}_{\OLMathNumeral{0}}\Proves\lforall[{\OLId{latin-ordinary}{x}}][!A({\OLId{latin-ordinary}{x}})]$؛ ثم ترفع الرتابة النتيجة إلى~${\OLId{greek-variable}{Gamma}}$.
\end{proof}

\begin{prop}
\ollabel{prop:provability-quantifiers}
لكل حد مغلق~${\OLId{latin-ordinary}{t}}$ تصح العبارتان الآتيتان:
\begin{tagenumerate}{prvEx,prvAll}
\tagitem{prvEx}{$!A({\OLId{latin-ordinary}{t}}) \Proves \lexists[{\OLId{latin-ordinary}{x}}][!A({\OLId{latin-ordinary}{x}})]$.}{}

\tagitem{prvAll}{$\lforall[{\OLId{latin-ordinary}{x}}][!A({\OLId{latin-ordinary}{x}})] \Proves !A({\OLId{latin-ordinary}{t}})$.}{}
\end{tagenumerate}
\end{prop}

\begin{proof}
\begin{tagenumerate}{prvEx,prvAll}
\tagitem{prvEx}{نأخذ \olref[qua]{ax:q2} ثم نطبق مبرهنة الاستنباط.}{}
\tagitem{prvAll}{ونأخذ \olref[qua]{ax:q1} ثم نطبق مبرهنة الاستنباط.}{}
\end{tagenumerate}
\end{proof}

\end{document}
